Un nombre complexe est ici, littéralement, un couple de réels : la forme algébrique $a + ib$
n'est donc pas une définition mais le premier théorème du chapitre, avec son unicité — celle
qui autorise l'identification des parties réelles et imaginaires, geste constant du calcul.

Le chapitre suit ensuite le programme : conjugué, module, argument, forme trigonométrique et
exponentielle, formules d'Euler et de Moivre, équation du second degré à discriminant négatif,
puis l'interprétation géométrique — distance, angle, alignement, orthogonalité, cercle — et
l'écriture complexe des transformations.

Un énoncé mérite qu'on s'y arrête, parce que sa forme n'est pas celle qu'on écrirait
naïvement. L'argument d'un produit est la somme des arguments \emph{modulo} $2\pi$ : avec un
argument choisi dans $]-\pi\,;\pi]$, l'égalité entre réels est fausse, et il suffit de deux
nombres d'argument $3\pi/4$ pour le voir. L'énoncé est donc écrit dans le groupe des angles,
ce qui est la traduction exacte du \og{} $[2\pi]$ \fg{} du lycée.
